\documentclass[11pt]{amsart}
\usepackage[T1]{fontenc}
\usepackage{lmodern,microtype,amsmath,amssymb,amsthm,booktabs}
\usepackage[margin=1in]{geometry}
\usepackage[hidelinks]{hyperref}
\newtheorem{theorem}{Theorem}[section]
\newtheorem{proposition}[theorem]{Proposition}
\newtheorem{lemma}[theorem]{Lemma}
\theoremstyle{remark}\newtheorem{remark}[theorem]{Remark}
\newcommand{\Z}{\mathbb Z}
\newcommand{\Q}{\mathbb Q}
\newcommand{\R}{\mathbb R}
\newcommand{\C}{\mathbb C}
\newcommand{\N}{\mathbb N}
\newcommand{\PP}{\mathcal P}
\DeclareMathOperator{\supp}{supp}
\title{Finite Collatz histories through support and Hurwitz zero clusters}
\author{Kokuno Yumeto}
\date{13 September 2026. Research-program note.}
\begin{document}
\begin{abstract}
Finite Collatz itineraries admit a polynomial affine lift that retains the
positions of all odd steps. We give its inverse, its residue realization,
and explicit fibres of several forgetful maps. We then apply the Split
Zero programme's compact Hurwitz reconstruction to exact odd-return residue
coordinates. On a known packet of $K$ words, the full local zero-cluster
jet through order $K-1$ recovers every probability weight and every joint
arithmetic predicate on the words. The inverse is explicit, works at a
nontrivial Riemann zeta zero of any multiplicity, and has a sharp worst-case
jet order. A three-word calculation recovers integral-cycle candidate mass
from the first two moments. These are finite information-preserving
constructions; they do not supply an orbit-descent estimate.
\end{abstract}
\maketitle

\section{The question and its source}
The Split Zero construction distinguishes an absent component from a
present component whose amplitude is zero. Its Collatz F-series note
\cite{support} already represents affine branches on support-homogeneous
coordinates. Its later Hurwitz reconstruction \cite{hurwitz} recovers a
compact probability measure from a complete local zero cluster. The
question here is what these constructions retain when their inputs are
actual finite Collatz histories.

There are two steps. First, history must be attached to faithful
coordinates: a label beside a collapsed scalar cannot restore an arbitrary
lost word. Second, the spectral family must be specified. We use a Hurwitz
deformation with an explicit inverse, rather than the eigenvalues of the
two-by-two branch product. The latter have a computable loss of information.
Classical parity coding supplies the dynamical coordinates; neither the
affine composition identity nor finite parity coding is claimed as new.
The contribution of this note is their explicit finite transport through
the cited reconstruction, with synchronized arithmetic tests retained.

\section{A faithful polynomial lift of a finite history}
Let $q$ be an indeterminate and $R=\Z[1/2][q]$. For $d\in\{0,1\}$ put
\[
 T_d^{(q)}(x)=\frac{q^d x+d}{2}.
\]
At $q=3$, selecting $d$ to be the parity of the current integer gives
the shortened Collatz map. A word $w=(d_0,\ldots,d_{N-1})$ is read in
chronological order; $T_{d_0}$ acts first. Write
\[
 m=\sum_{j=0}^{N-1}d_j,\qquad s_j=\sum_{i<j}d_i,\qquad
 P_w(q)=\sum_{j=0}^{N-1}d_j2^j q^{m-1-s_j},
\]
where a summand with $d_j=0$ is omitted, so no negative exponent occurs.

\begin{proposition}[Lift, inverse, and realization]\label{prop:lift}
For every word $w$,
\begin{equation}\label{eq:lift}
 T_{d_{N-1}}^{(q)}\circ\cdots\circ T_{d_0}^{(q)}(x)
 =\frac{q^m x+P_w(q)}{2^N}.
\end{equation}
The map $w\mapsto(N,P_w)$ is injective. Its image consists of $(N,0)$
and pairs
\[
 \left(N,\sum_{k=0}^{m-1}2^{j_k}q^{m-1-k}\right),
 \qquad 0\le j_0<\cdots<j_{m-1}<N.
\]
The inverse places ones precisely at the $j_k$ and zeros elsewhere.
A positive integer $n$ realizes $w$ as its first $N$ shortened-map
parities exactly when
\begin{equation}\label{eq:parity-residue}
 3^m n+P_w(3)\equiv0\pmod{2^N}.
\end{equation}
Thus each word occurs on one residue class modulo $2^N$.
\end{proposition}
\begin{proof}
Appending $d$ replaces $(m,P)$ by $(m+d,q^dP+d2^N)$, which proves
\eqref{eq:lift} and the coefficient formula by induction. If the one
positions are $j_0<\cdots<j_{m-1}$, the coefficient of $q^{m-1-k}$
is exactly $2^{j_k}$. This proves the image and inverse, including the
all-zero word. The length is essential: terminal zero bits contribute
no new coefficient to $P$.

For realization use induction on $N$. An actual orbit satisfies
\eqref{eq:parity-residue} by \eqref{eq:lift}. Conversely suppose the
congruence holds. If $d_0=0$, then $P_w(3)$ is even, so $n$ is even;
write $n=2y$ and $P_w(3)=2P_{w'}(3)$ for the remaining word $w'$.
Division by two yields its congruence at $y$. If $d_0=1$, then
$P_w(3)=3^{m-1}+2P_{w'}(3)$ is odd, so $n$ is odd; put
$y=(3n+1)/2$. The same division gives
$3^{m-1}y+P_{w'}(3)\equiv0\pmod{2^{N-1}}$.
Induction proves every requested parity. The coefficient $3^m$ is
invertible modulo $2^N$, giving the unique residue. Positive representatives
exist and remain positive along these actual iterates.
\end{proof}

The source F-series uses the opposite matrix order, with $d_0$ acting
last on zero:
\[
 X_w(q)=T_{d_0}^{(q)}\circ\cdots\circ T_{d_{N-1}}^{(q)}(0)
 =\sum_{j=0}^{N-1}\frac{d_jq^{s_j}}{2^{j+1}}.
\]
If $w^{\mathrm{rev}}$ reverses the word, then
\[
 X_{w^{\mathrm{rev}}}(q)=2^{-N}P_w(q).
\]
Indeed, reversing a one at $j$ puts it at $N-1-j$, with $m-1-s_j$
ones before it. This is the exact map between the two composition
conventions, not an identification of input binary digits with orbit parities.

\subsection{Support, specialization, and what each forgets}
For the source semiring $G(R)=R\sqcup\{\tau\}$, $\tau$ is the additive
identity and multiplicative absorber, while $e=0_R$ is a supported zero.
For $r,s\in R$, their addition and multiplication are those of $R$.
The reflection $p:G(R)\to R$ fixes $R$ and sends $\tau$ to $0$.
Its zero fibre is exactly $\{\tau,e\}$; every other fibre is a singleton.
The Boolean support character has
\[
 \chi(\tau)=0,\qquad \chi(r)=1\quad(r\in R),
\]
including $\chi(e)=1$. Hence this character detects presence, not
arithmetic nonvanishing or parity. For a nonzero coefficient ring, the
source branch matrices are
\[
 M_0=\begin{pmatrix}1/2&e\\\tau&1\end{pmatrix},\qquad
 M_1=\begin{pmatrix}q/2&1/2\\\tau&1\end{pmatrix}.
\]
Both have the same Boolean support matrix. Their full polynomial products
retain more than this support projection.

To make a history mask explicit, set $A_N=R^{\{0,\ldots,N-1\}}$ and
$\epsilon_w=(d_0,\ldots,d_{N-1})\in A_N$. It is a central idempotent.
Its $j$th coordinate is the parity bit itself, so $(N,\epsilon_w)$
recovers the word. This mask is additional retained data; applying $\chi$
to the two branch matrices does not produce it.

The evaluation homomorphism $R\to\Z[1/2]$, $q\mapsto3$, has kernel
$(q-3)$ by division by the monic polynomial $q-3$.
For the source F-series, take the five-bit words with one positions
$\{0,4\}$ and $\{2,3,4\}$ respectively. Then
\[
 X_u(q)=\frac12+\frac q{32},\qquad
 X_v(q)=\frac18+\frac q{16}+\frac{q^2}{32},
\]
and
\[
 X_v(q)-X_u(q)=\frac{(q-3)(q+4)}{32},\qquad
 X_u(3)=X_v(3)=\frac{19}{32}.
\]
These are distinct histories in one specialization fibre.
Under $G(R)\to G(R/(q-3))$ their nonzero difference becomes the
supported zero $e$, not $\tau$. Keeping that support label records
that a contribution was present and killed; recovering the contribution
itself requires its preimage or the retained polynomial coefficients.

The reflected forward product in \eqref{eq:lift} has matrix
\[
 \begin{pmatrix}q^m/2^N&P_w(q)/2^N\\0&1\end{pmatrix}.
\]
Its eigenvalue multiset is $\{q^m/2^N,1\}$, independent of the order
of the bits. At $N=2,m=1$, $w=(0,1)$ has translation $1/2$ and
$w=(1,0)$ has translation $1/4$, with the same eigenvalues. This
computes the loss for this spectral map; it does not preclude a different,
faithful spectral construction.

\section{Exact odd-return coordinates}
Fix integers $m\ge1$ and $A\ge m$, and let
\[
 W_{m,A}=\{(a_1,\ldots,a_m)\in\N_{>0}^m:\textstyle\sum a_i=A\},
 \qquad K=|W_{m,A}|=\binom{A-1}{m-1}.
\]
These are labelled words, not rotation classes. Put $S_j=a_1+\cdots+a_j$,
$S_0=0$, and
\[
 B_w=\sum_{j=0}^{m-1}3^{m-1-j}2^{S_j},\qquad D=2^A-3^m.
\]
The odd-return branch composition is $(3^m x+B_w)/2^A$.
Let $r_w$ be the unique representative $1\le r_w<2^{A+1}$ satisfying
\begin{equation}\label{eq:odd-residue}
 3^m r_w+B_w\equiv2^A\pmod{2^{A+1}}.
\end{equation}

\begin{lemma}\label{lem:odd}
The positive odd integers with first $m$ exact odd-return exponents $w$
are precisely $r_w+2^{A+1}\Z$, intersected with the positive integers.
The residues $r_w$ are pairwise distinct on $W_{m,A}$.
\end{lemma}
\begin{proof}
For one letter, \eqref{eq:odd-residue} is exactly
$\nu_2(3x+1)=a_1$. For more letters write
$B_w=3^{m-1}+2^{a_1}B_{w'}$. The congruence first implies
$2^{a_1}\mid3x+1$. Put $y=(3x+1)/2^{a_1}$ and divide by $2^{a_1}$.
The resulting congruence is
\[
 3^{m-1}y+B_{w'}\equiv2^{A-a_1}\pmod{2^{A-a_1+1}}.
\]
Since $A-a_1\ge1$ and $B_{w'}$ is odd, it forces $y$ odd. The
valuation was therefore exact, and induction applies to the suffix.
The converse follows by reversing this calculation. Multiplication by
$3^m$ is invertible modulo $2^{A+1}$, proving uniqueness and oddness
of $r_w$. If two words shared a residue, its positive representative
would have two different exact exponent lists, impossible for a function.
\end{proof}

The rational periodic candidate is $x_w=B_w/D$. The denominator is
nonzero, odd, and coprime to three. Rotation $\sigma$ satisfies
\[
 3B_w+D=2^{a_1}B_{\sigma w}.
\]
Because all $B$ are odd, this gives exactly the exponent $a_1$ in
the rationals with odd denominator. The candidates are positive when
$D>0$ and integral precisely when $D\mid B_w$. In the latter case
the identity and coprimality propagate integrality around the word.
These observations justify the predicate used below without mistaking
all rational candidates for integer cycles.

\section{A zero-cluster inverse for packet distributions}
Let $\Delta(W)$ be the probability simplex on the known word set $W=W_{m,A}$.
For $\lambda\in\Delta(W)$ define, in the original integer residue coordinate,
\begin{equation}\label{eq:family}
 \mu_\lambda=\sum_{w\in W}\lambda_w\delta_{r_w},\qquad
 M_j=\sum_w\lambda_w r_w^j,\qquad
 Z_\lambda(s,t)=\sum_w\lambda_w\zeta(s,1+tr_w).
\end{equation}
Here $M_0=1$. Take $L=2^{A+1}-1$ and $|t|L<1$. The parameter $t$
is a complex deformation parameter, not Collatz time. Hurwitz Taylor
expansion \cite{dlmf} gives
\begin{equation}\label{eq:hurwitz}
 Z_\lambda(s,t)=\sum_{j\ge0}M_j b_j(s)t^j,\qquad
 b_j(s)=\frac{(-1)^j}{j!}(s)_j\zeta(s+j).
\end{equation}
On a compact subset of $0<\Re s<1$, the factor $(s)_j/j!$ is bounded
by a polynomial in $j$: multiply $1+(s-1)/k$ for $1\le k\le j$ and
use the harmonic-sum bound on its absolute value. For $j\ge1$ the zeta
factor is uniformly bounded by its absolutely convergent Dirichlet series.
Since $|M_j|\le L^j$, the series is normally convergent locally for
$|t|L<1$. This establishes joint holomorphicity in the neighbourhoods used
below. At any nontrivial zero $\rho$, every $b_j(\rho)$, $j\ge1$, is
nonzero: the rising factorial has no zero factor, and $\Re(\rho+j)>1$,
where the absolutely convergent Euler product makes zeta nonvanishing.

\begin{lemma}[Cluster-to-moment inverse, after the Split Zero reconstruction]
\label{lem:cluster}
Fix a nontrivial zero $\rho$ of exact multiplicity $r$. Let $Q_\lambda(s,t)$
be the monic degree-$r$ polynomial in $s-\rho$ whose roots are the entire
nearby zero cluster of $Z_\lambda(\cdot,t)$, with multiplicities. Then its
$t$-jet through order $J$ and $(M_1,\ldots,M_J)$ determine one another.
\end{lemma}
\begin{proof}
Choose a disc wholly inside the critical strip that contains only $\rho$
among the zeros at $t=0$, with no boundary zero. Uniform convergence and
Rouche's theorem give precisely $r$ interior zeros for small $t$.
The contour integrals of
$(s-\rho)^k\partial_s Z_\lambda/Z_\lambda$, divided by $2\pi i$, give
the power sums of their displacements, for $1\le k\le r$. Newton's
identities construct holomorphic elementary symmetric functions and hence
$Q_\lambda$, without choosing individual branches. Write
\[
 Q_\lambda=(s-\rho)^r+\sum_{j\ge1}Q_j(s)t^j,\qquad
 Z_\lambda=U_\lambda Q_\lambda,\qquad
 U_\lambda=\sum_{j\ge0}U_j(s)t^j.
\]
The quotient is jointly holomorphic: on an enclosing contour it is
holomorphic in $t$, and its Cauchy integral gives its extension in $s$.
For each fixed $t$ the two divisors coincide, so all apparent singularities
are removable. It is a unit, and $U_0=\zeta(s)/(s-\rho)^r$ is known.
Coefficient comparison yields the successive inverse formulas
\begin{align}
 M_j&=\frac{\sum_{k=1}^j U_{j-k}(\rho)Q_k(\rho)}{b_j(\rho)},
 \label{eq:inverseM}\\
 U_j(s)&=\frac{M_jb_j(s)-\sum_{k=1}^j U_{j-k}(s)Q_k(s)}{(s-\rho)^r}.
 \label{eq:inverseU}
\end{align}
The second numerator vanishes to order at least $r$, by the actual analytic
factorization; its value after division at $\rho$ is its $r$th derivative
divided by $r!$. Starting with $U_0$, these are explicit triangular
recurrences. They are inverses on cluster jets arising from
\eqref{eq:family}; arbitrary formal jets are not asserted to be admissible.

Conversely, agreement of moments through $J$ gives agreement of $Z$ to
order $t^{J+1}$ on the contour, also after $s$-differentiation. The
denominators there stay nonzero, so the contour power sums, Newton
coefficients, and monic polynomials agree to the same order. This proves
the converse. The argument reproduces the portion of \cite{hurwitz},
H10--H13, needed here, and requires neither RH nor simplicity.
\end{proof}

\begin{theorem}[Finite spectral reconstruction]\label{thm:main}
The map from $\Delta(W_{m,A})$ to actual zero-cluster jets of order
$K-1$ in \eqref{eq:family} is injective, with singleton fibres on its
image. Its inverse is \eqref{eq:inverseM}--\eqref{eq:inverseU} followed by
\begin{equation}\label{eq:lagrange}
 \ell_w(y)=\prod_{v\ne w}\frac{y-r_v}{r_w-r_v}
 =\sum_{j=0}^{K-1}\ell_{w,j}y^j,\qquad
 \lambda_w=\sum_{j=0}^{K-1}\ell_{w,j}M_j.
\end{equation}
For $K\ge2$, order $K-2$ does not suffice for all probability vectors.
For $K=1$, the unique weight is already known from $M_0=1$.
\end{theorem}
\begin{proof}
Distinct nodes were proved in Lemma~\ref{lem:odd}. Thus every denominator
in \eqref{eq:lagrange} is nonzero and $\ell_w(r_v)$ is the Kronecker
delta. Integrating gives the inverse; the preceding lemma gives its moments.

For sharpness let $h_w=\prod_{v\ne w}(r_w-r_v)^{-1}$. Comparing the
leading coefficient in Lagrange interpolation of $y^j$ gives
\[
 \sum_w h_w r_w^j=0\quad(0\le j\le K-2),\qquad
 \sum_w h_w r_w^{K-1}=1.
\]
The vectors $\lambda_w^\pm=1/K\pm\varepsilon h_w$ are distinct strictly
positive probabilities whenever
$0<\varepsilon<1/(K\max_w|h_w|)$. They have equal moments through
$K-2$, hence identical cluster jets of that order, and distinct next
moments. This proves the claimed lower bound.
\end{proof}

This is an exact inverse, not a numerical conditioning bound. In particular,
the shrinking guaranteed $t$-disc as $A$ grows and the Lagrange coefficients
remain present in the construction.

\section{Arithmetic predicates and an explicit three-word calculation}
For any function $f:W\to\R$, define its interpolating polynomial
\[
 H_f(y)=\sum_w f(w)\ell_w(y)=\sum_{j=0}^{K-1}h_jy^j.
\]
Then the exact expectation is
\begin{equation}\label{eq:predicate}
 \sum_w\lambda_w f(w)=\sum_{j=0}^{K-1}h_jM_j.
\end{equation}
For odd moduli $q_1,\ldots,q_b$, choose
$f(w)=\prod_i\mathbf 1_{B_w\equiv c_i\pmod{q_i}}$ to recover the
joint residue mass. The product is evaluated on a single word, so no
independence of the marginal tests is introduced. Pairwise coprimality is
needed only to identify the tuple with one Chinese-remainder residue,
not to define or recover this joint mass.

For $D>0$, the choice $f(w)=\mathbf 1_{D\mid B_w}$ recovers the mass
of integral periodic candidates. With uniform weights, multiplying this
mass by $K$ gives the number of qualifying labelled words, including
repetitions. It is not a count of primitive cycles. A particular predicate
can require fewer moments if its interpolating polynomial has lower degree;
the sharpness theorem concerns the entire probability vector.

For $(m,A)=(2,4)$, $D=7$ and the complete data are
\[
\begin{array}{c|r|r|c}
 w&B_w&r_w&7\mid B_w\\\hline
 (1,3)&5&19&\text{no}\\
 (2,2)&7&1&\text{yes}\\
 (3,1)&11&29&\text{no}
\end{array}
\]
Thus
\begin{equation}\label{eq:example}
 \lambda_{(2,2)}=\frac{M_2-48M_1+551}{504}.
\end{equation}
Indeed, the relevant interpolant is
$(y-19)(y-29)/((1-19)(1-29))$. Uniform weights give $M_1=49/3$,
$M_2=401$, and mass $1/3$. The qualifying word repeats the known odd
fixed point $1$; the example is an exact transfer calculation.

To write \eqref{eq:example} directly from a simple zero trajectory,
let $\rho+c_1t+c_2t^2+\cdots$ be that trajectory. Then
\begin{align*}
 M_1&=\frac{\zeta'(\rho)c_1}{\rho\zeta(\rho+1)},\\
 M_2&=-\frac{2}{\rho(\rho+1)\zeta(\rho+2)}
 \left(\zeta'(\rho)c_2+\frac{\zeta''(\rho)}2c_1^2
 -M_1\bigl(\zeta(\rho+1)+\rho\zeta'(\rho+1)\bigr)c_1\right).
\end{align*}
The general-multiplicity formulas \eqref{eq:inverseM}--\eqref{eq:inverseU}
give the same mass without a simplicity assumption.

\subsection{Maps of histories and the role of zero weights}
For a map of finite labelled sets $\pi:W\to V$ (for example, deleting
the last exponent), the induced probability map is
\[
 (\pi_*\lambda)_v=\sum_{w:\pi(w)=v}\lambda_w.
\]
Its linear extension has kernel
$\{h:\sum_{\pi(w)=v}h_w=0\text{ for all }v\}$; on probability
vectors each nonempty fibre is exactly the set with these prescribed
fibre sums. This computes the histories lost by taking prefixes.
The two spectral encodings commute with this map after decoding: recover
$\lambda$, apply these sums, and construct the target Hurwitz family
using any specified distinct coordinates on $V$. This defines a map on
the actual images, not on arbitrary analytic germs.

To transport values without decoding each component separately, interpolate
$w\mapsto f(\pi(w))$ in \eqref{eq:predicate}. This preserves precisely
the same expectations. If a set $S\subseteq W$ also records histories
retained despite zero weight, the input is $(S,\lambda)$ with
$\supp\lambda\subseteq S$. The spectral family depends only on
$\lambda$. Its fibre in this enlarged input consists of the possible
$S\supseteq\supp\lambda$. Retaining $S$ separately is therefore necessary
to distinguish supported zero weight from absence, exactly as in the
Split Zero construction. It is not information present in the measure alone.

\section{Mathematical use and verification}
The construction supplies faithful finite spectral coordinates for the
same joint arithmetic tests used in the Collatz residue workbench. It
also gives explicit witnesses when a chosen projection loses history.
The Hurwitz family is built from the history weights themselves; recovering
them from its zeros does not discover an unknown distribution of infinite
positive-integer trajectories. A quantitative use for descent would need
additional estimates about those trajectories. No such estimate is proved
here, and no equivalence between RH and Collatz is inferred from using a
zeta deformation.

The accompanying standard-library certificate checks finite parity-word
compositions, exact residue realizations, polynomial reconstruction,
the specialization collision, Lagrange inverses, and the displayed packet
calculation using integer and rational arithmetic. The general assertions
are proved above. The analytic zero-cluster argument is not Lean-certified
or established by sampling numerical zeta zeros. The source records and
independent review identify which part of the preceding Split Zero
programme was used.

\begin{thebibliography}{9}
\bibitem{support} \emph{Support-Idempotent Homogenization for a Split-Zero
Collatz F-Series Frame}, local technical note, April 2026, Sections 1,
4, 5, and 7. The exact source file and SHA256 are in the accompanying
source manifest. The present note uses the displayed finite branch
matrices and prefix formula; it does not rely on all claims in that note.
\bibitem{hurwitz} Split Zero research programme,
\emph{Historical Hurwitz reconstruction and the complete Laurent jet bridge},
equations H1--H13, especially H10--H13. Public source snapshot,
13 September 2026:
\href{https://github.com/KokunoYumeto/zeta-function-research-reader/blob/16fc4dbb817b82023c6126e636f1c6df28d4cecd/workbenches/splitzero-tandem/continuations/20260913-toda/tex/historical_hurwitz_jets.tex}{\texttt{zeta-function-research-reader}, commit \texttt{16fc4dbb817b}}.
\bibitem{companion} Kokuno Yumeto,
\emph{Weighted Odd-Step Paths, Survivor Coding, Arithmetic Groupoids, and
Toric Period Morphisms: A Research Companion for the Collatz Map},
Sections 2 and 4, working edition, 4 September 2026.
\href{https://www.overleaf.com/read/vgsnwwgbrnsm\#2eec34}{Collatz workbench}.
\bibitem{dlmf} NIST Digital Library of Mathematical Functions,
\href{https://dlmf.nist.gov/25.11.E10}{equation 25.11.10}, Hurwitz zeta
Taylor expansion. Accessed 13 September 2026.
\end{thebibliography}
\end{document}
